\documentclass[12pt]{article}
\usepackage[utf8]{inputenc}
\usepackage[T1]{fontenc}
\usepackage[margin=1in]{geometry}
\usepackage{setspace}
\usepackage{xcolor}
\usepackage[colorlinks=true, urlcolor=blue!60!black, linkcolor=blue!60!black]{hyperref}
\setlength{\parindent}{0em}
\setlength{\parskip}{0.6em}
\renewcommand{\baselinestretch}{1.0}
\begin{document}
\begin{center}
{\LARGE \textit{United States Steel Corp. v. Multistate Tax Commission}, 434 U.S. 452 (1978)\par}
\vspace{0.5em}
{\large Case Brief --- Two-Page Edition\par}
\end{center}

\section{Memory Jogger}
The Multistate Tax Compact needed no congressional consent: it gave member states no power they lacked individually and did not threaten federal supremacy (\href{https://supreme.justia.com/cases/federal/us/434/452/}{Justia}).

\section{Detailed Case Facts}
When federal uniformity bills stalled, state tax administrators drafted the Multistate Tax Compact (1966), creating a Commission whose regulations were advisory and which could audit corporations at a member state's request. States kept control of rates and assessments and could withdraw at will. Congress never passed any of twelve consent bills.

\section{Procedural History}
United States Steel and other multistate taxpayers sued; a three-judge district court granted summary judgment for the Commission; the Court affirmed on direct appeal.

\section{Judicial Votes}
Affirmed, 7--2. Powell wrote for Burger, Brennan, Stewart, Marshall, Rehnquist, and Stevens; White, joined by Blackmun, dissented.

\section{Holding}
Under \textit{Virginia v.\ Tennessee}, consent is required only for agreements tending to increase state political power in a way that may encroach on federal supremacy. This Compact did not; the Commerce Clause and Fourteenth Amendment claims also failed.

\section{Analysis of Opinions}
\textbf{Majority:} A literal reading would require consent for every trivial agreement. The test concerns potential impact on federal supremacy; multilateral membership and an agency do not suffice. The Compact delegated no sovereign power, allowed withdrawal, and used methods each state could adopt alone, while Congress could still preempt. \textbf{White:} the Court demanded proof of actual harm; the Commission lobbied against federal bills, and the Clause must bar some joint actions states could take individually.

\section{Examples: Future Applications}
\textbf{Same-Side:} an advisory multistate commission on AI with voluntary audits needs no consent. \textbf{Opposite-Side:} a maritime compact that levies tariffs on foreign vessels and bars withdrawal requires consent. \textbf{Fence-Sitter:} the National Popular Vote Interstate Compact.

\section{Critique}
Critics say the test nullifies the text of ``any Agreement or Compact,'' ignores harms to nonmember states, and treats a permanent multistate agency as no different from states acting alone.

\section{Key Quotations}
\begin{quote}
``But the test is whether the Compact enhances state power \textit{quoad} the National Government. This pact does not purport to authorize the member States to exercise any powers they could not exercise in its absence. Nor is there any delegation of sovereign power to the Commission; each State retains complete freedom to adopt or reject the rules and regulations of the Commission.''
\end{quote}
\noindent --- Justice Powell, Opinion of the Court, 434 U.S. at 473.

\end{document}
